% TOP_PROBLEM: {"id": 10000014, "title": "Grid-or-tree embeddings in superlinear Cayley graphs", "category_id": 3, "category": {"id": 3, "name": "graph_theory", "display_name": "Graph Theory", "description": "Problems involving graphs, networks, and their properties.", "slug": "graph-theory", "order_index": 3, "created_at": "2026-07-31T15:26:25.671Z"}, "status": "open", "problem_number": "AMR-099-0014", "set_id": 13}
% FIRST_POSTED: 2026-10-01
% ENTRY_KIND: statement_only
% UnsolvedMath ID 10000014; source catalogue code AMR-099-0014
% !TeX program = lualatex
% Definition and sources only; this file is not an LLM solution attempt.
\documentclass[11pt,a4paper]{article}
\usepackage[margin=25mm]{geometry}
\usepackage{fontspec}
\setmainfont{lmroman10-regular.otf}[BoldFont=lmroman10-bold.otf,
 ItalicFont=lmroman10-italic.otf,BoldItalicFont=lmroman10-bolditalic.otf]
\usepackage{amsmath,amssymb,mathtools}
\usepackage{unicode-math}
\setmathfont{latinmodern-math.otf}
\AtBeginDocument{\let\setminus\smallsetminus}
\usepackage{xurl,hyperref}
\hypersetup{colorlinks=true,urlcolor=blue}
\setcounter{secnumdepth}{0}
\setlength{\parindent}{0pt}
\setlength{\parskip}{5pt}
\setlength{\emergencystretch}{3em}
\begin{document}
\section*{Grid-or-tree embeddings in superlinear Cayley graphs}\label{10000014}
{\small UnsolvedMath ID 10000014 · AMR-099-0014 · Graph Theory · AMR Open Problem Lists}
\subsection{Definitions and mathematical statement}
Let $\Gamma$ be a finitely generated infinite group and $S=S^{-1}$ a finite generating set not containing its identity. Let $G=\operatorname{Cay}(\Gamma,S)$ be the undirected graph with edges $\{g,gs\}$ for $g\in\Gamma$, $s\in S$. Give $G$ its graph metric. Suppose it has superlinear volume growth:
\[
\frac{|B_G(1,n)|}{n}\longrightarrow\infty.
\]
Let $Q$ be the nearest-neighbor square lattice on $\mathbb Z^2$, and let $T$ be the infinite rooted binary tree, with each vertex having two children.

Must at least one of $Q,T$ admit a rough isometric embedding in $G$? Here, as in the source's coarse-geometric terminology, this means a map $f:V(H)\to V(G)$ for $H=Q$ or $H=T$, and constants $A\ge1$, $B\ge0$, such that for all $x,y\in V(H)$,
\[
A^{-1}d_H(x,y)-B\le d_G(f(x),f(y))\le A d_H(x,y)+B.
\]
This is a quasi-isometric embedding; coarse surjectivity onto $G$ is not required. The constants may depend on the group and its generating set. The map need not carry each edge to one edge, so the formulation is a metric embedding problem rather than a demand for a literal lattice or binary-tree subgraph.

The two possible targets have different growth: a binary-tree embedding forces exponential growth in the ambient graph, while the lattice alternative also allows intermediate or polynomial growth. The question concerns every finite generating set of every superlinear Cayley graph.
\subsection{Short English statement}
Does every Cayley graph growing faster than a line contain a coarse metric copy of a square grid or a binary tree?
\subsection{Sources}
\begin{itemize}
\item[S1] ULAM.AI and the UnsolvedMath Contributors, supplied problems.json, record 10000014 (AMR-099-0014), AMR Open Problem Lists. Catalogue identity and source transcription; CC BY 4.0.
\url{https://huggingface.co/datasets/ulamai/UnsolvedMath}.
\item[S2] Benjamini - Coarse Geometry and Randomness (Saint-Flour notes), Open problem 1.80, PDF page 17. Original source indexed by AMR-099-0014.
\url{https://www.wisdom.weizmann.ac.il/~itai/stflouraug24.pdf#page=17}.
\item[S3] I. Benjamini, Coarse Geometry and Randomness, primary Saint-Flour notes; the numbered question and its surrounding definitions.
\url{https://arquivo.pt/noFrame/replay/20201231041548id_/http://www.wisdom.weizmann.ac.il/~itai/stflouraug24.pdf}.
\end{itemize}
\end{document}
